\chapter{Modele nasycone i jednorodne}
\label{ch:mod-saturation}
\section{Motywacja i nasycenie}
Twierdzenie o realizacji typów w rozszerzeniach elementarnych
(\ref{thm:mod-type-elementary-extension}) nie gwarantuje
realizacji w wyjściowym modelu. Nasycenie jest warunkiem
zapewniającym, że przy odpowiednio małym zbiorze parametrów
nie trzeba wychodzić poza model, aby znaleźć realizację typu.
Jednorodność natomiast pozwala przenosić elementy mające
ten sam typ. Obie własności będą potrzebne przy badaniu
struktury modeli teorii zupełnej.

Pracujemy w ZFC, w ustalonej sygnaturze $\mathcal L$.
Uniwersa struktur $\mathfrak M,\mathfrak N$ oznaczamy
odpowiednio przez $M,N$. Parametr $\kappa$ oznacza moc nieskończoną, utożsamianą
z odpowiadającą jej początkową liczbą porządkową.
Przez $\kappa^+$ oznaczamy najmniejszą moc kardynalną
ściśle większą od $\kappa$.

\begin{definicja}[Nasycenie]
\label{def:mod-saturation}
Strukturę $\mathfrak M$ nazywamy $\kappa$-nasyconą, jeśli
dla każdego $A\subseteq M$ o $|A|<\kappa$ każdy typ
zupełny $p\in S_1(A)$ jest realizowany w $\mathfrak M$.
Nieskończoną strukturę $\mathfrak M$ nazywamy nasyconą,
jeżeli jest $|M|$-nasycona.
\end{definicja}

Warunek dotyczy typu nad $A$ w sensie definicji
\ref{def:mod-type-over-parameters}; nie chodzi jedynie
o zgodność z dowolnie wybraną teorią struktury bez parametrów.
Jeśli $\lambda\leq\kappa$ są mocami nieskończonymi,
to $\kappa$-nasycenie implikuje $\lambda$-nasycenie,
gdyż każdy zbiór parametrów mocy mniejszej od $\lambda$
ma moc mniejszą od $\kappa$.

\begin{definicja}[Odwzorowanie częściowo elementarne]
\label{def:mod-partial-elementary}
Dla $A\subseteq M$ funkcję $f:A\to N$ nazywamy
częściowo elementarną z $\mathfrak M$ do $\mathfrak N$,
jeśli dla każdej formuły $\varphi(\bar x)$ i każdej
krotki $\bar a$ z $A$ odpowiedniej długości zachodzi
\[
\mathfrak M\models\varphi(\bar a)
\quad\Longleftrightarrow\quad
\mathfrak N\models\varphi(f(\bar a)).
\]
Dopuszczamy pustą krotkę, a więc także zdania.
Zapis $f(\bar a)$ oznacza krotkę obrazów współrzędnych.
W szczególności istnienie dowolnego odwzorowania częściowo
elementarnego implikuje $\mathfrak M\equiv\mathfrak N$,
a puste odwzorowanie jest częściowo elementarne dokładnie
wtedy, gdy struktury są elementarnie równoważne.
\end{definicja}

\begin{lemat}[Podstawowe własności]
\label{lem:mod-partial-elementary}
Odwzorowanie częściowo elementarne jest różnowartościowe;
jego odwrotność określona na obrazie jest częściowo
elementarna. Ograniczenia i złożenia takich odwzorowań
są częściowo elementarne (gdy złożenie jest określone
na całej rozważanej dziedzinie).
Suma rosnącego, niepustego łańcucha odwzorowań częściowo
elementarnych między tymi samymi strukturami jest
odwzorowaniem częściowo elementarnym.
\end{lemat}
\begin{proof}
Różnowartościowość wynika z zastosowania równoważności
z definicji do formuły $x_1=x_2$. Dla odwrotności odwracamy
tę równoważność, a dla złożenia stosujemy ją kolejno
do obu funkcji. Ograniczenie zachowuje wszystkie wymagane
równoważności. Suma rosnącego łańcucha funkcji jest funkcją,
ponieważ wartości na wspólnych częściach dziedzin są zgodne.
Każda skończona krotka z dziedziny sumy należy do dziedziny
pewnej funkcji w łańcuchu: wybieramy największy spośród
skończenie wielu potrzebnych etapów. W przypadku pustej
krotki wybieramy dowolny etap. Elementarność na tym etapie
daje równoważność dla sumy.
\end{proof}

\section{Rozszerzanie odwzorowań częściowo elementarnych}
\begin{lemat}[Przedłużenie o jeden element]
\label{lem:mod-saturation-forth}
Niech $\mathfrak N$ będzie $\kappa$-nasycona,
$f:A\to N$ częściowo elementarna z $\mathfrak M$ do
$\mathfrak N$ oraz $|A|<\kappa$. Dla każdego $a\in M$
istnieje częściowo elementarne przedłużenie
$g:A\cup\{a\}\to N$ funkcji $f$.
\end{lemat}
\begin{proof}
Jeżeli $a\in A$, wystarczy wziąć $g=f$. Załóżmy zatem,
że $a\notin A$. Oznaczmy $B=f[A]$; przez różnowartościowość
$|B|=|A|<\kappa$. Przenosimy typ elementu $a$ nad $A$
na zbiór $B$, zastępując w każdej formule nazwę parametru
$c_u$ nazwą $c_{f(u)}$. Otrzymujemy zbiór formuł
\[
\begin{aligned}
q(x)=\{\varphi(x,f(\bar u)):\;&
       \mathfrak M\models\varphi(a,\bar u),\\
       &\bar u\text{ jest skończoną krotką z }A\}.
\end{aligned}
\]
Precyzyjnie jest to obraz typu
$\operatorname{tp}^{\mathfrak M}(a/A)$ przez przemianowanie
nowych symboli stałych, ustalające wszystkie symbole
pierwotnego języka $\mathcal L$.
W tym zapisie $\varphi(x,\bar u)$ oznacza wynik podstawienia
nazw parametrów za wyróżnione wolne zmienne formuły
języka $\mathcal L$, jak w rozdziale \ref{ch:mod-types}.
Każdą formułę z nazwami można tak zapisać, zastępując
skończenie wiele nazw świeżymi zmiennymi.
Twierdzenie \ref{thm:sem-substitution} uzasadnia zgodność
spełniania przy obu zapisach.
Aby sprawdzić skończoną spełnialność $q$ w $\mathfrak N$,
weźmy skończenie wiele jego formuł i zbierzmy wszystkie
ich parametry w jednej krotce $\bar u$ z $A$.
W $\mathfrak M$ prawdziwe jest zdanie z parametrami
$\exists x\,\bigwedge_i\varphi_i(x,\bar u)$, gdyż
świadkiem jest $a$. Częściowa elementarność $f$ przenosi
tę formułę egzystencjalną do $\mathfrak N$ przy
parametrach $f(\bar u)$. Zatem skończony fragment
jest realizowany w $\mathfrak N$.
Z twierdzenia \ref{thm:mod-type-elementary-extension}
wynika, że $q$ jest typem częściowym nad $B$.
Jest też zupełny: przemianowanie komutuje z negacją,
a każda formuła z parametrami z $B$
ma jednoznaczne tłumaczenie przez odwrotność $f$ na $A$,
a typ $a$ rozstrzyga tę formułę lub jej negację.

Nasycenie daje element $b\in N$ realizujący $q$. Ponieważ dla
każdego $u\in A$ formuła $x\ne c_{f(u)}$ należy do $q$,
mamy $b\notin B$. Kładziemy $g(a)=b$ i $g|_A=f$.
Każda formuła przy krotce z $A\cup\{a\}$, w której
występuje $a$, daje formułę w jednej wyróżnionej zmiennej
z parametrami z $A$: wszystkie wolne wystąpienia zmiennych
wartościowanych przez $a$ zastępujemy tą samą świeżą zmienną,
a pozostałe zmienne wolne nazwami odpowiednich elementów $A$.
Zmienne związane przemianowujemy w razie potrzeby,
aby uniknąć przechwycenia zmiennej wolnej
(por. \ref{thm:sem-alpha-renaming-xlv}).
Definicja $q$ i jego realizacja zapewniają zgodność
prawdziwości tej formuły, także gdy jest fałszywa,
gdyż wtedy do $q$ należy jej negacja.
Zgodność tego podstawienia z wartościowaniem wynika
z twierdzenia \ref{thm:sem-substitution}. Dla krotek z $A$
korzystamy z elementarności $f$.
Otrzymujemy żądane przedłużenie.
\end{proof}

\begin{twierdzenie}[Typy skończonych krotek]
\label{thm:mod-saturation-finite-tuples}
Jeżeli $\mathfrak M$ jest $\kappa$-nasycona, to dla
każdego $A\subseteq M$ o $|A|<\kappa$ i każdego $n\geq1$
realizuje każdy typ częściowy w $n$ zmiennych nad $A$.
W szczególności definicja \ref{def:mod-saturation} jest
równoważna wersji wymagającej realizacji wszystkich
typów zupełnych skończonych krotek.
\end{twierdzenie}
\begin{proof}
Niech $p$ będzie typem częściowym w $n$ zmiennych nad $A$.
Na mocy twierdzenia \ref{thm:mod-type-elementary-extension}
istnieje $\mathfrak N\succeq\mathfrak M$ i krotka
$(a_1,\ldots,a_n)\in N^n$ realizująca $p$.
Identyczność na $A$ jest częściowo elementarna z
$\mathfrak N$ do $\mathfrak M$, ponieważ
$\mathfrak M\preccurlyeq\mathfrak N$. Stosujemy
lemat \ref{lem:mod-saturation-forth} kolejno do
$a_1,\ldots,a_n$, za każdym razem rozszerzając dziedzinę
o co najwyżej jeden element. Wszystkie dziedziny mają
moc mniejszą od $\kappa$, ponieważ $\kappa$ jest
nieskończona, a do $A$ dodajemy skończenie wiele elementów.
Końcowe odwzorowanie $f$ przenosi realizację $p$
na krotkę $(f(a_1),\ldots,f(a_n))$ w $\mathfrak M$
i ustala parametry z $A$. Implikacja odwrotna wynika
z przypadku $n=1$ i typów zupełnych.
\end{proof}

\begin{przyklad}[Granica nasycenia]
\label{ex:mod-saturation-size}
Jeżeli $\mathfrak M$ jest nieskończona i $|M|=\mu$,
to nie jest $\mu^+$-nasycona. Istotnie zbiór formuł
$\{x\ne c_a:a\in M\}$ jest skończenie realizowalny
w $\mathfrak M$, ponieważ można ominąć dowolny
skończony podzbiór $M$. Jest typem częściowym nad $M$
na mocy twierdzenia \ref{thm:mod-type-elementary-extension},
ale nie jest realizowany w $\mathfrak M$.
Ponieważ $|M|<\mu^+$, przeczyłoby to twierdzeniu
\ref{thm:mod-saturation-finite-tuples}.
\end{przyklad}

\begin{przyklad}[Język samej równości]
\label{ex:mod-saturation-equality}
Każda nieskończona struktura $\mathfrak M$ języka samej
równości, o mocy $\mu$, jest nasycona. Dla $A\subseteq M$,
$|A|<\mu$, z przykładu \ref{ex:mod-type-pure-equality}
wiemy, że każdy typ z $S_1(A)$ jest typem pewnego
elementu $A$ albo jedynym typem zawierającym wszystkie
nierówności $x\ne c_a$ dla $a\in A$.
Pierwsze typy są realizowane przez odpowiednie elementy,
a ostatni przez dowolny element $M\setminus A$,
który istnieje z nierówności mocy.
\end{przyklad}

\begin{przyklad}[Standardowa arytmetyka nie jest $\aleph_0$-nasycona]
\label{ex:mod-saturation-arithmetic}
W języku arytmetyki z $0$, następnikiem $S$ i porządkiem
rozważamy typ częściowy nad pustym zbiorem parametrów
\[
p(x)=\{S^n(0)<x:n\in\N\}.
\]
Zapis $S^n(0)$ oznacza term otrzymany przez $n$-krotne
zastosowanie następnika. Każdy skończony fragment
realizuje odpowiednio duża liczba naturalna, ale żaden
element standardowej arytmetyki nie realizuje całego $p$.
Z twierdzenia \ref{thm:mod-type-elementary-extension}
$p$ jest typem częściowym nad zbiorem pustym.
Twierdzenie \ref{thm:mod-saturation-finite-tuples}
wyklucza $\aleph_0$-nasycenie.
\end{przyklad}

\section{Łańcuchy elementarne i konstrukcja rozszerzeń}
\begin{lemat}[Suma dowolnego łańcucha elementarnego]
\label{lem:mod-saturation-chain}
Niech $\delta$ będzie niezerową liczbą porządkową,
a $(\mathfrak M_\alpha)_{\alpha<\delta}$ rodziną
struktur, dla której $\mathfrak M_\alpha\preccurlyeq
\mathfrak M_\beta$ dla $\alpha\leq\beta<\delta$.
Na sumie $M=\bigcup_{\alpha<\delta}M_\alpha$ istnieje
jedyna struktura $\mathfrak M$ zgodna z interpretacjami
na wszystkich etapach i
$\mathfrak M_\alpha\preccurlyeq\mathfrak M$
dla każdego $\alpha<\delta$.
\end{lemat}
\begin{proof}
Jest to uogólnienie lematu \ref{lem:ls-elementary-chain}.
Każda skończona krotka z sumy należy do pewnego etapu,
a interpretacje funkcji, stałych i relacji na etapach
są zgodne, co definiuje strukturę na sumie. Każda krotka
argumentów funkcji należy do pewnego etapu, więc także
jej wartość należy do sumy; wybór etapu nie wpływa na
wynik. Relacje definiujemy analogicznie, a stałe mają
tę samą wartość na każdym etapie. Konstrukcja jest
jednoznaczna, ponieważ obejmuje wszystkie krotki.
Indukcję po formułach prowadzimy jednocześnie dla
wszystkich etapów. Dla atomów działa definicja
struktury na sumie, dla spójników hipoteza indukcyjna.
Jeśli formuła egzystencjalna z parametrami z $M_\alpha$
jest prawdziwa w $\mathfrak M_\alpha$, jej świadek
należy do $M_\alpha$ i jest świadkiem w sumie na mocy
hipotezy indukcyjnej. Jeśli jest prawdziwa w sumie,
jej świadek należy do pewnego $M_\beta$ z $\beta\geq\alpha$
(bierzemy maksimum indeksów). Hipoteza indukcyjna
przenosi formułę ze świadkiem do $M_\beta$,
a elementarność $\mathfrak M_\alpha$ w $\mathfrak M_\beta$
przenosi zdanie egzystencjalne z powrotem do $M_\alpha$.
Kwantyfikator ogólny wynika z negacji i egzystencjalnego.
\end{proof}

\begin{lemat}[Jednoczesna realizacja rodziny typów]
\label{lem:mod-saturation-many-types}
Niech $(p_i)_{i\in I}$ będzie rodziną indeksowaną zbiorem $I$
typów częściowych w jednej zmiennej nad podzbiorami
$A_i\subseteq M$. Istnieje rozszerzenie elementarne
$\mathfrak N\succeq\mathfrak M$ realizujące każdy $p_i$.
Można je wybrać tak, aby
\[
|N|\leq\max(\aleph_0,|\mathcal L|,|M|,|I|).
\]
\end{lemat}
\begin{proof}
Dodajemy do języka z nazwami wszystkich elementów $M$
parami różne świeże stałe $d_i$ dla $i\in I$.
Rozważamy teorię
\[
\Gamma=\operatorname{ED}(\mathfrak M)
\cup\bigcup_{i\in I}\{\varphi(d_i):\varphi(x)\in p_i\}.
\]
W zapisie $\varphi(d_i)$ podstawiamy $d_i$ wyłącznie
za wolne wystąpienia zmiennej $x$ i zachowujemy nazwy
parametrów z $A_i$. Diagram pochodzi z definicji
\ref{def:ls-elementary-diagram}. Każdy skończony fragment $\Gamma$ wymaga
realizacji skończenie wielu formuł z każdego spośród
skończenie wielu typów. Na mocy twierdzenia
\ref{thm:mod-type-elementary-extension} dla każdego
z tych typów znajdujemy odpowiedni element $M$
i interpretujemy przez niego właściwą stałą $d_i$.
Nie wymagamy różnych realizacji, więc wybory nie
kolidują. Pozostałe nowe stałe interpretujemy przez
ustalony element niepustego $M$. Diagram elementarny
jest spełniony w $\mathfrak M_M$.
Z twierdzenia o zwartości \ref{thm:compactness} otrzymujemy
model $\mathfrak D\models\Gamma$,
a lemat \ref{lem:ls-diagram-embedding} daje zanurzenie
elementarne $j:\mathfrak M\to\mathfrak D|_{\mathcal L}$.
Stosujemy dolne twierdzenie Löwenheima--Skolema
\ref{thm:downward-ls} do reduktu, z parametrami
\[
j[M]\cup\{d_i^{\mathfrak D}:i\in I\}.
\]
Otrzymujemy podstrukturę elementarną $\mathfrak C$ o mocy
nie większej od wskazanej granicy, zawierającą wszystkie
realizacje. Ponieważ $j[\mathfrak M]$ i $\mathfrak C$ są
podstrukturami elementarnymi tego samego reduktu
i $j[M]\subseteq C$, mamy $j[\mathfrak M]\preccurlyeq\mathfrak C$:
dla każdej krotki z $j[M]$ prawdziwość w obu strukturach
jest równoważna prawdziwości w redukcie. Ta sama
elementarność zachowuje formuły typów przy świadkach
$d_i^{\mathfrak D}$ i parametrach z $j[M]$.
Przenosimy strukturę $\mathfrak C$ izomorficznie,
utożsamiając $j(a)$ z $a$
dla $a\in M$, jak w dowodzie twierdzenia
\ref{thm:mod-type-elementary-extension}. Elementarność
i prawdziwość formuł z parametrami zachowują się przy
tym przeniesieniu. Otrzymujemy żądane $\mathfrak N$.
\end{proof}

\begin{twierdzenie}[Istnienie rozszerzeń nasyconych]
\label{thm:mod-saturation-extension}
Dla każdej struktury $\mathfrak M$ i każdej mocy
nieskończonej $\kappa$ istnieje $\kappa$-nasycone
rozszerzenie elementarne $\mathfrak N\succeq\mathfrak M$.
Jeżeli $\lambda$ jest mocą nieskończoną,
$\lambda\geq\max(|M|,2^{\max(\aleph_0,|\mathcal L|,\kappa)})$ i
$\lambda^\kappa=\lambda$, to można uzyskać
$|N|\leq\lambda$.
\end{twierdzenie}
\begin{proof}
Wybierzmy $\lambda$ o podanych własnościach. Zawsze istnieje
taka moc: dla $\mu=\max(\aleph_0,|M|,|\mathcal L|,\kappa)$
możemy wziąć $\lambda=2^\mu$, ponieważ
$(2^\mu)^\kappa=2^{\mu\cdot\kappa}=2^\mu$.
Pierwsza równość wynika z bijekcji między rodzinami
podzbiorów $\mu$ indeksowanymi przez $\kappa$ a podzbiorami
$\kappa\times\mu$; druga z $\kappa\leq\mu$ i nieskończoności
$\mu$. Z definicji $\mu$ spełniony jest również warunek
dolnego ograniczenia na $\lambda$.

Budujemy rosnący łańcuch elementarny
$(\mathfrak M_\alpha)_{\alpha<\kappa^+}$,
zaczynając od $\mathfrak M_0=\mathfrak M$ i zachowując
$|M_\alpha|\leq\lambda$. Na etapie następnikowym
rozważamy wszystkie pary $(A,p)$, gdzie
$A\subseteq M_\alpha$, $|A|<\kappa$ i $p\in S_1(A)$.
Liczba zbiorów $A$ jest nie większa od $\lambda^\kappa$:
każdy taki zbiór można zakodować przez funkcję z $\kappa$
do $M_\alpha$ z dodatkowym symbolem pustym, której
niepuste wartości tworzą ten zbiór.
Z lematu \ref{lem:ls-formula-count} liczba formuł języka
$\mathcal L(A)$ jest co najwyżej
$\max(\aleph_0,|\mathcal L|,\kappa)$.
Typ jest podzbiorem zbioru formuł, więc liczba typów
nad każdym $A$ jest co najwyżej
$2^{\max(\aleph_0,|\mathcal L|,\kappa)}\leq\lambda$.
Łącznie liczba par $(A,p)$ nie przekracza
$\lambda\cdot\lambda=\lambda$.
Z lematu \ref{lem:mod-saturation-many-types} uzyskujemy
$\mathfrak M_{\alpha+1}$ realizującą wszystkie te typy
i o mocy co najwyżej $\lambda$.

Na niezerowych etapach granicznych bierzemy sumy łańcuchów,
stosując lemat \ref{lem:mod-saturation-chain}. Przechodniość
elementarności i ten lemat zapewniają, że dla każdych
$\alpha\leq\beta<\kappa^+$ zachodzi
$\mathfrak M_\alpha\preccurlyeq\mathfrak M_\beta$. Ponieważ
$\lambda^\kappa=\lambda$, z twierdzenia Cantora mamy
$\lambda\geq 2^\kappa\geq\kappa^+$ (iniekcja zbioru
funkcji $\kappa\to\{0,1\}$ w zbiór funkcji $\kappa\to\lambda$
daje pierwszą nierówność, a twierdzenie Cantora drugą).
Każda suma najwyżej
$\kappa^+$ zbiorów mocy co najwyżej $\lambda$ ma moc
co najwyżej $\lambda$.
W celu prowadzenia rekurencji na zbiorze, a nie na
klasach właściwych, po każdym etapie następnikowym
przenosimy nowe elementy do osobnego zbioru mocy
$\lambda$, rozłącznego z $M$ i z analogicznymi zbiorami
dla innych etapów. Wówczas wszystkie wybory modeli
mogą być dokonywane spośród struktur na podzbiorach
jednego ustalonego zbioru; zbiór wszystkich struktur języka $\mathcal L$ na podzbiorach
tego ustalonego zbioru dobrze porządkujemy i na każdym
etapie wybieramy pierwszy dopuszczalny model, którego
nowe elementy należą do zbioru zarezerwowanego dla
tego etapu. Interpretacje na granicach są jednoznaczne.

Niech $\mathfrak N$ będzie sumą całego łańcucha.
Jeżeli $A\subseteq N$ i $|A|<\kappa$, wybieramy dla
każdego $a\in A$ indeks etapu zawierającego $a$.
Wszystkie te indeksy są ograniczone przez pewien
$\alpha<\kappa^+$: w przeciwnym razie $\kappa^+$ byłaby
sumą mniej niż $\kappa$ liczb porządkowych, z których
każda ma moc co najwyżej $\kappa$ (można brać ich następniki).
Taka suma miałaby moc co najwyżej $\kappa\cdot\kappa=\kappa$,
co jest sprzecznością. Dla $A=\varnothing$ bierzemy $\alpha=0$.
Zatem $A\subseteq M_\alpha$.
Z elementarności $\mathfrak M_\alpha\preccurlyeq\mathfrak N$
przestrzeń $S_1(A)$ w obu strukturach jest identyczna
(por. definicja \ref{def:mod-type-over-parameters}
i następująca po niej obserwacja).
Każdy typ z $S_1(A)$ został zatem zrealizowany w
$\mathfrak M_{\alpha+1}$ (mamy $\alpha+1<\kappa^+$,
gdyż nieskończona początkowa liczba porządkowa jest graniczna)
i jego realizacja jest zachowana
w $\mathfrak N$. To daje $\kappa$-nasycenie.
\end{proof}

\begin{twierdzenie}[Nasycenie dokładnie zadanej mocy]
\label{thm:mod-saturated-fixed-cardinality}
Niech $\lambda$ będzie nieprzeliczalną mocą regularną,
$|\mathcal L|<\lambda$ oraz niech
$\lambda^\nu\leq\lambda$ dla każdej mocy $\nu<\lambda$.
Jeżeli $\mathfrak M$ jest nieskończona i $|M|\leq\lambda$,
to istnieje nasycone rozszerzenie elementarne
$\mathfrak N\succeq\mathfrak M$ o mocy dokładnie $\lambda$.
\end{twierdzenie}
\begin{proof}
Z górnego twierdzenia Löwenheima--Skolema
\ref{thm:upward-ls} wybieramy rozszerzenie elementarne
$\mathfrak M_0\succeq\mathfrak M$ o mocy $\lambda$.
Modyfikujemy konstrukcję z twierdzenia
\ref{thm:mod-saturation-extension}: budujemy tym razem
łańcuch długości $\lambda$, realizując na każdym etapie
następnikowym wszystkie typy zupełne w jednej zmiennej
nad podzbiorami poprzedniego etapu mocy mniejszej od
$\lambda$. Sprawdzimy, że wystarcza moc $\lambda$.
Jeśli $|M_\alpha|=\lambda$, to liczba takich podzbiorów
nie przekracza
\[
\sum_{\nu<\lambda,\ \nu\ \text{moc}} \lambda^\nu
\leq\lambda\cdot\lambda=\lambda.
\]
Istotnie każdy podzbiór mocy $\nu$ jest obrazem pewnej
funkcji $\nu\to M_\alpha$, a liczba możliwych mocy
$\nu<\lambda$ nie przekracza $\lambda$.
Dla każdego takiego podzbioru $A$ liczba formuł
w języku $\mathcal L(A)$ jest nie większa od
$\theta=\max(\aleph_0,|\mathcal L|,|A|)<\lambda$
(lemat \ref{lem:ls-formula-count}).
Liczba typów nad $A$ nie przekracza
$2^\theta\leq\lambda^\theta\leq\lambda$.
Lemat \ref{lem:mod-saturation-many-types} daje zatem
rozszerzenie mocy co najwyżej $\lambda$; ma ono moc
co najmniej $\lambda$, ponieważ zawiera $M_0$.
Na niezerowych etapach granicznych i na końcu bierzemy
sumy łańcuchów zgodnie z lematem \ref{lem:mod-saturation-chain}.
Ich moc także jest równa $\lambda$, gdyż sumujemy
co najwyżej $\lambda$ zbiorów mocy $\lambda$.
Rekurencję prowadzimy na ustalonym zbiorze możliwych
struktur, z rozłącznymi rezerwami mocy $\lambda$ dla
każdego etapu, jak w dowodzie
twierdzenia \ref{thm:mod-saturation-extension}.

Niech $A\subseteq N$ ma moc mniejszą od $\lambda$.
Z regularności $\lambda$, czyli
$\operatorname{cf}(\lambda)=\lambda$, zbiór indeksów
etapów, na których pojawiają się elementy $A$, jest
ograniczony w $\lambda$. Zatem $A\subseteq M_\alpha$
dla pewnego $\alpha<\lambda$ (dla pustego $A$ bierzemy
$\alpha=0$). Typy nad $A$ są takie same w
$\mathfrak M_\alpha$ i $\mathfrak N$, ponieważ
$\mathfrak M_\alpha\preccurlyeq\mathfrak N$.
Wszystkie zrealizowano na etapie $\alpha+1<\lambda$.
W konsekwencji $\mathfrak N$ jest $\lambda$-nasycona
i $|N|=\lambda$, czyli jest nasycona.
\end{proof}

\begin{uwaga}
Założenia poprzedniego twierdzenia są warunkami
wystarczającymi, nie koniecznymi. Przykład
\ref{ex:mod-saturation-equality} daje modele nasycone
każdej mocy nieskończonej bez tych założeń.
Nie wnioskujemy stąd o istnieniu modelu nasyconego
dowolnej teorii i dowolnej zadanej mocy.
\end{uwaga}

\begin{przyklad}[Teoria przeliczalna bez przeliczalnego modelu nasyconego]
\label{ex:mod-saturation-cardinality-obstruction}
Niech $\mathcal L$ będzie językiem z równością,
którego jedynymi symbolami pozalogicznymi są
jednoargumentowe symbole relacyjne $P_n$ dla $n\in\N$.
Jako uniwersum $M$ bierzemy zbiór wszystkich skończonych
podzbiorów $\N$, a $P_n^{\mathfrak M}(F)$ niech zachodzi
dokładnie wtedy, gdy $n\in F$. Zbiór $M$ jest przeliczalnie
nieskończony: jest sumą po $k\in\N$ zbiorów podzbiorów
$k$-elementowych, z których każdy jest co najwyżej
przeliczalny jako obraz podzbioru $\N^k$, a wszystkie
zbiory jednoelementowe tworzą rodzinę nieskończoną.
Ustalmy teorię zupełną $T=\operatorname{Th}(\mathfrak M)$.
Dla każdego $D\subseteq\N$ rozważmy zbiór formuł
\[
p_D(x)=
\{P_n(x):n\in D\}\cup\{\neg P_n(x):n\in\N\setminus D\}.
\]
Każdy skończony podzbiór $p_D$ jest realizowany w $\mathfrak M$
przez skończony zbiór indeksów występujących w nim
formuł pozytywnych: nie zawiera on żadnego indeksu
występującego w tym fragmencie z negacją.
Z lematu \ref{lem:mod-type-compactness} $p_D$ jest typem
częściowym nad $T$, czyli także nad zbiorem pustym
w dowolnym modelu $T$ (por. definicja
\ref{def:mod-type-over-parameters}).

Jeśli $D\ne E$, to istnieje indeks $n$ należący do dokładnie
jednego z tych zbiorów. Po ewentualnej zamianie $D$ i $E$
mamy $n\in D\setminus E$. Typy $p_D$ i $p_E$ zawierają
odpowiednio formułę $P_n(x)$ i jej negację, zatem nie
mogą mieć wspólnej realizacji. Gdyby model $\mathfrak N$
teorii $T$ był $\aleph_0$-nasycony, to na mocy
twierdzenia \ref{thm:mod-saturation-finite-tuples}
realizowałby wszystkie $p_D$. Wybierając po jednej
realizacji każdego typu (w ZFC), otrzymalibyśmy
iniekcję $\mathcal P(\N)\to N$. Stąd
$|N|\geq 2^{\aleph_0}>\aleph_0$ na mocy twierdzenia
Cantora. W szczególności teoria $T$ nie ma
przeliczalnego modelu nasyconego, mimo że ma
przeliczalny model $\mathfrak M$.
\end{przyklad}

\section{Jednorodność i metoda tam i z powrotem}
\begin{definicja}[Jednorodność i silna jednorodność]
\label{def:mod-homogeneity}
Struktura $\mathfrak M$ jest $\kappa$-jednorodna,
jeżeli każde częściowo elementarne $f:A\to M$
o $|A|<\kappa$ można przedłużyć częściowo elementarnie
na $A\cup\{a\}$ dla dowolnego $a\in M$.
Jest silnie $\kappa$-jednorodna, jeśli każde takie $f$
można przedłużyć do automorfizmu $\mathfrak M$.
\end{definicja}

Terminologia jednorodności bywa różna; w tym skrypcie
stosujemy dokładnie powyższe rozróżnienie.
Lemat \ref{lem:mod-saturation-forth} pokazuje, że
$\kappa$-nasycenie implikuje $\kappa$-jednorodność.
Silna $\kappa$-jednorodność implikuje $\kappa$-jednorodność,
ponieważ automorfizm można ograniczyć do $A\cup\{a\}$.

\begin{przyklad}[Silna jednorodność bez nasycenia]
\label{ex:mod-homogeneity-not-saturation}
Standardowa arytmetyka w języku zawierającym $0$ i
następnik $S$ jest silnie $\kappa$-jednorodna dla każdej
mocy nieskończonej $\kappa$. Istotnie, jeżeli $f:A\to\N$
jest częściowo elementarne, to dla każdego $n\in A$
formuła $x=S^n(0)$ spełniona przy $n$ musi być spełniona
przy $f(n)$. Ma ona w standardowej arytmetyce dokładnie
jedną realizację, zatem $f(n)=n$. Funkcja $f$ jest więc
ograniczeniem identyczności, która jest automorfizmem.
Jeżeli język zawiera również porządek, przykład
\ref{ex:mod-saturation-arithmetic} pokazuje, że struktura
nie jest nawet $\aleph_0$-nasycona. Silna jednorodność nie
implikuje więc nasycenia.
\end{przyklad}

\begin{twierdzenie}[Tam i z powrotem dla modeli nasyconych]
\label{thm:mod-saturation-back-forth}
Niech $\mathfrak M,\mathfrak N$ będą strukturami
$\kappa$-nasyconymi o mocy $\kappa$.
Każde częściowo elementarne odwzorowanie $f:A\to N$
z $\mathfrak M$ do $\mathfrak N$, dla $|A|<\kappa$,
przedłuża się do izomorfizmu $\mathfrak M\to\mathfrak N$.
\end{twierdzenie}
\begin{proof}
Ustalamy wyliczenia $M=\{a_\alpha:\alpha<\kappa\}$
i $N=\{b_\alpha:\alpha<\kappa\}$.
Budujemy rosnący łańcuch częściowo elementarnych
odwzorowań $f_\alpha$ dla $\alpha<\kappa$,
zaczynając od $f_0=f$.

Ustalamy dobre porządki $M$ i $N$, aby przy każdym
przedłużeniu wybierać pierwszy dopuszczalny obraz.
Na etapie $\alpha+1$ stosujemy lemat
\ref{lem:mod-saturation-forth} do $f_\alpha$ i $a_\alpha$,
otrzymując przedłużenie $g$ zawierające $a_\alpha$
w dziedzinie (krok tam). Następnie do $g^{-1}$,
częściowo elementarnego z $\mathfrak N$ do $\mathfrak M$
na mocy lematu \ref{lem:mod-partial-elementary},
stosujemy ten sam lemat i element $b_\alpha$.
Odwracając otrzymane przedłużenie, dostajemy
$f_{\alpha+1}$, którego obraz zawiera $b_\alpha$
(krok z powrotem). Na etapach granicznych bierzemy
sumy wcześniejszych funkcji.

Przed każdym krokiem dziedzina i obraz mają moc
mniejszą od $\kappa$: do początkowej dziedziny
dodaliśmy najwyżej dwa elementy na każdy wcześniejszy
etap, więc dla $\alpha<\kappa$ moc nie przekracza
$|A|+2\cdot|\alpha|+2<\kappa$.
Nie wymaga to regularności $\kappa$.
Łańcuch jest częściowo elementarny dzięki lematowi
\ref{lem:mod-partial-elementary}.

Suma $F=\bigcup_{\alpha<\kappa}f_\alpha$ jest
częściowo elementarna i różnowartościowa, a każdy
$a_\alpha$ należy do jej dziedziny i każdy $b_\alpha$
do jej obrazu. Zatem jest bijekcją $M\to N$.
Zachowanie formuł atomowych daje zachowanie i odbijanie
relacji. Dla każdej arności używamy odpowiedniej liczby
parami różnych zmiennych (w tym świeżej zmiennej $y$
dla wartości funkcji). Dla symbolu funkcyjnego $h$ stosujemy formułę
$h(\bar x)=y$ do krotki argumentów i jej wartości
w $\mathfrak M$; otrzymujemy
$F(h^{\mathfrak M}(\bar a))=h^{\mathfrak N}(F(\bar a))$; dla stałej analogicznie formułę $x=c$.
Stąd $F$ jest izomorfizmem.
\end{proof}

Rysunek \ref{fig:mod-back-forth} przedstawia etap,
w którym oba wskazane elementy wymagają dodania.

\begin{figure}[htbp]
\centering
\begin{tikzpicture}[>=Stealth,scale=0.92]
\draw[rounded corners,fill=blue!5,draw=blue!50!black]
(0,0) rectangle (4,3);
\draw[rounded corners,fill=green!5,draw=green!45!black]
(6,0) rectangle (10,3);
\node at (2,3.35) {$M$};
\node at (8,3.35) {$N$};
\fill (2,2.4) circle (2pt) node[left] {$u$};
\fill (8,2.4) circle (2pt) node[right] {$f_\alpha(u)$};
\draw[->,thick] (2.15,2.4)--(7.85,2.4);
\fill[blue!75!black] (2,1.65) circle (2pt) node[left] {$a_\alpha$};
\fill[blue!75!black] (8,1.65) circle (2pt) node[right] {$g(a_\alpha)$};
\draw[->,blue!75!black,thick] (2.15,1.65)--(7.85,1.65)
node[midway,above] {tam};
\fill[green!45!black] (2,0.65) circle (2pt) node[left] {$v$};
\fill[green!45!black] (8,0.65) circle (2pt) node[right] {$b_\alpha$};
\draw[<-,green!45!black,thick] (2.15,0.65)--(7.85,0.65)
node[midway,above] {z powrotem};
\end{tikzpicture}

\caption{Schemat jednego etapu metody tam i z powrotem:
najpierw obejmujemy nowy element dziedziny, następnie
nowy element obrazu.}
\label{fig:mod-back-forth}
\end{figure}

\begin{wniosek}[Jednoznaczność i silna jednorodność]
\label{cor:mod-saturated-uniqueness}
Dwa modele nasycone tej samej teorii zupełnej o tej samej
nieskończonej mocy są izomorficzne.
Każdy model nasycony mocy $\kappa$ jest silnie
$\kappa$-jednorodny.
\end{wniosek}
\begin{proof}
Dwa modele teorii zupełnej są elementarnie równoważne,
więc pusta funkcja jest częściowo elementarna
(definicja \ref{def:mod-partial-elementary}).
Stosujemy twierdzenie \ref{thm:mod-saturation-back-forth}.
Dla silnej jednorodności stosujemy je z $\mathfrak N=\mathfrak M$
i dowolnym częściowo elementarnym odwzorowaniem o
dziedzinie mocy mniejszej od $\kappa$.
\end{proof}

\begin{wniosek}[Typy a orbity automorfizmów]
\label{cor:mod-types-orbits}
Niech $\mathfrak M$ będzie modelem nasyconym mocy $\kappa$,
$A\subseteq M$, $|A|<\kappa$, a $\bar a,\bar b\in M^n$.
Wówczas
$\operatorname{tp}^{\mathfrak M}(\bar a/A)=
\operatorname{tp}^{\mathfrak M}(\bar b/A)$
wtedy i tylko wtedy, gdy istnieje automorfizm
$\sigma$ struktury $\mathfrak M$, ustalający punktowo $A$
i spełniający $\sigma(\bar a)=\bar b$.
\end{wniosek}
\begin{proof}
Automorfizm ustalający $A$ zachowuje prawdziwość
formuł z tymi parametrami przez twierdzenie
\ref{thm:model-iso-invariance}.
Odwrotnie, równość typów zapewnia zgodność wszystkich
równości między współrzędnymi oraz między współrzędnymi
i parametrami z $A$. Funkcja ustalająca $A$ i wysyłająca
$a_i$ na $b_i$ jest więc dobrze określona i
różnowartościowa. Każda formuła z krotką elementów
jej dziedziny daje formułę dotyczącą $\bar a$ z parametrami
z $A$: elementom należącym do $A$ przypisujemy ich nazwy,
a pozostałym wybieramy odpowiadające im współrzędne
$\bar a$. Zgodność równości zapewnia, że te same
wybory współrzędnych opisują obrazy przez $\bar b$.
Podstawienia wykonujemy za wolne zmienne, po ewentualnym
przemianowaniu związanych, zgodnie z twierdzeniami
\ref{thm:sem-substitution} i \ref{thm:sem-alpha-renaming-xlv}. Równość typów zapewnia elementarność częściową.
Dziedzina ma moc mniejszą od $\kappa$, więc wniosek
\ref{cor:mod-saturated-uniqueness} daje automorfizm.
\end{proof}

\section{Uniwersalność i ograniczenia konwencji modelu-monstra}
\begin{twierdzenie}[Uniwersalność modelu nasyconego]
\label{thm:mod-saturated-universal}
Jeżeli $\mathfrak N$ jest $\kappa$-nasycona,
$\mathfrak M\equiv\mathfrak N$ i $|M|\leq\kappa$,
to istnieje zanurzenie elementarne $\mathfrak M$
w $\mathfrak N$. Co więcej, każde odwzorowanie częściowo
elementarne $f:A\to N$ z $\mathfrak M$ do $\mathfrak N$,
dla $|A|<\kappa$, przedłuża się do takiego zanurzenia.
\end{twierdzenie}
\begin{proof}
Wystarczy udowodnić drugie stwierdzenie, gdyż pusta
funkcja jest częściowo elementarna przez
$\mathfrak M\equiv\mathfrak N$. Zaczynamy od $f$.
Ustalamy dobry porządek $N$ i na każdym etapie
wybieramy pierwszy dopuszczalny obraz, jeśli trzeba
dodać element spoza aktualnej dziedziny.
Ustalamy wyliczenie
$M$ indeksowane liczbą porządkową $|M|$.
Na każdym etapie następnikowym zapewniamy, że kolejny
element wyliczenia należy do dziedziny, stosując lemat
\ref{lem:mod-saturation-forth} (jeśli już należy, nic
nie zmieniamy),
a na niezerowym etapie granicznym bierzemy sumę.
Definiujemy również odwzorowanie na końcowym etapie
$|M|$, także gdy jest on etapem następnikowym.
Przed krokiem o indeksie $\alpha<|M|\leq\kappa$
dziedzina ma moc co najwyżej $|A|+|\alpha|<\kappa$:
dodaliśmy najwyżej jeden element na każdy wcześniejszy
etap, a suma dwóch mocy mniejszych od nieskończonej
mocy $\kappa$ pozostaje mniejsza od $\kappa$.
Odwzorowanie na końcowym etapie przedłuża $f$,
jest określone na całym $M$ i częściowo elementarne. Jest zatem zanurzeniem elementarnym:
zachowuje stałe, funkcje i relacje przez formuły atomowe,
a wszystkie formuły przez definicję.
\end{proof}

\begin{uwaga}[Model-monster jako konwencja]
\label{rem:mod-monster-model}
W teorii modeli często ustala się jeden dostatecznie duży
model teorii zupełnej, w którym prowadzi się rozważania
o małych zbiorach parametrów i typach. Nazywa się go
modelem-monstrem. Dla ustalonej granicy $\kappa$ oczekuje
się co najmniej $\kappa$-nasycenia i silnej
$\kappa$-jednorodności; małość oznacza moc mniejszą
od $\kappa$. Jeżeli dysponujemy modelem nasyconym mocy
$\lambda\geq\kappa$, to wniosek
\ref{cor:mod-saturated-uniqueness} daje obie własności.
Samo twierdzenie \ref{thm:mod-saturation-extension}
zapewnia $\kappa$-nasycenie pewnego rozszerzenia, lecz
nie dowodzi istnienia modelu nasyconego każdej zadanej
mocy ani silnej jednorodności dowolnego modelu
$\kappa$-nasyconego. Jeżeli teoria ma nieskończony model
mocy co najwyżej $\lambda$, gdzie $\lambda\geq\kappa$
spełnia warunki twierdzenia
\ref{thm:mod-saturated-fixed-cardinality}, można skonstruować
model nasycony mocy $\lambda$ i wykorzystać go w tej roli.
W dalszych rozważaniach takie założenia będziemy
formułować jawnie.
\end{uwaga}
